%==============================================================================
\chapter{Generated Programs Examined in the Case Studies}
\label{app:generated_programs}
%==============================================================================

This appendix reproduces the complete selected programs discussed in
Section~\ref{sec:generated_program_cases}. The genetic algorithm was generated
by Qwen2.5-Coder-32B-Instruct and the evolution strategy by
DeepSeek-Coder-V2-Lite-Instruct.
The code is reproduced as selected, including the implementation weaknesses
discussed in Chapter~\ref{chap:experiments}. Line numbers refer to original
source lines; long lines wrap for presentation. The listings make the
relationship between selection, objective calls, state updates, and the
returned solution explicit.

\section{Genetic Champion with Observation Averaging}

This champion was generated by Qwen2.5-Coder-32B-Instruct for
Gallagher 101 peaks ($f_{21}$), $D=2$, and $\sigma=0.2$, using Explicit
mode and the Vectorization scaffold (experiment 1298, iteration 6).
In the current held-out benchmark, all 20 trials return finite solutions;
15 meet $\Delta f\leq10^{-8}$ and 19 meet $\Delta f\leq10^{-2}$.
Each trial uses 20,000 objective evaluations. The averaging mechanism,
diversity-based mutation rule, and historical-incumbent limitation are
examined in Section~\ref{subsec:generated_ga_case}.

\begin{lstlisting}[style=championprogram,caption={Complete generated \texttt{ImprovedNoisyGeneticAlgorithm} champion with five-observation averaging.},label={lst:generated_ga_complete}]
import numpy as np

class ImprovedNoisyGeneticAlgorithm:
    def __init__(self):
        self.population_size = 100
        self.mutation_rate = 0.2
        self.crossover_rate = 0.7
        self.elite_fraction = 0.1
        self.tournament_size = 5
        self.evaluations_per_candidate = 5  # Increase evaluations per candidate

    def __call__(self, problem, budget):
        lb = np.asarray(getattr(problem, 'lower_bound', [-5.0, -5.0]), dtype=float)
        ub = np.asarray(getattr(problem, 'upper_bound', [5.0, 5.0]), dtype=float)
        dim = int(getattr(problem, 'dim', len(lb) if hasattr(lb, '__len__') else 3))

        # Initialize population
        population = np.random.uniform(lb, ub, size=(self.population_size, dim))
        fitnesses = self.evaluate_population(population, problem)
        evaluations = self.population_size * self.evaluations_per_candidate

        # Sort initial population by fitness
        sorted_indices = np.argsort(fitnesses)
        population = population[sorted_indices]
        fitnesses = fitnesses[sorted_indices]

        best_x, best_y = population[0], fitnesses[0]

        while evaluations < budget:
            new_population = []

            # Elitism: carry over the top fraction of individuals
            elite_size = int(self.elite_fraction * self.population_size)
            new_population.extend(population[:elite_size])

            # Fill the rest of the population with crossover and mutation
            while len(new_population) < self.population_size:
                parent1 = self.tournament_selection(population, fitnesses)
                parent2 = self.tournament_selection(population, fitnesses)

                if np.random.rand() < self.crossover_rate:
                    child = self.crossover(parent1, parent2)
                else:
                    child = parent1.copy()

                # Mutate the child
                child = self.mutate(child, lb, ub)

                new_population.append(child)

            # Evaluate the new population
            new_fitnesses = self.evaluate_population(new_population, problem)
            evaluations += len(new_population) * self.evaluations_per_candidate

            # Sort by fitness and replace the old population
            sorted_indices = np.argsort(new_fitnesses)
            population = np.array(new_population)[sorted_indices]
            fitnesses = new_fitnesses[sorted_indices]

            # Update best_x, best_y if improvement found
            if fitnesses[0] < best_y:
                best_x, best_y = population[0], fitnesses[0]

            # Adaptively adjust mutation rate based on population diversity
            diversity = np.mean(np.std(population, axis=0))
            self.mutation_rate = 0.1 + 0.3 * (diversity / (ub - lb)).mean()

        return best_x, float(best_y)

    def evaluate_population(self, population, problem):
        fitnesses = []
        for candidate in population:
            evaluations = [problem(candidate) for _ in range(self.evaluations_per_candidate)]
            fitnesses.append(np.mean(evaluations))
        return np.array(fitnesses)

    def crossover(self, parent1, parent2):
        alpha = np.random.rand(len(parent1))
        child = alpha * parent1 + (1 - alpha) * parent2
        return child

    def mutate(self, individual, lb, ub):
        mutation_mask = np.random.rand(len(individual)) < self.mutation_rate
        noise = np.random.normal(0, 0.5, size=len(individual))
        mutated_individual = np.where(mutation_mask, individual + noise * (ub - lb), individual)
        return np.clip(mutated_individual, lb, ub)

    def tournament_selection(self, population, fitnesses):
        selected_indices = np.random.choice(range(len(population)), self.tournament_size, replace=False)
        best_index = selected_indices[np.argmin(fitnesses[selected_indices])]
        return population[best_index]
\end{lstlisting}

\section{Instance-Sensitive Evolution-Strategy Champion}

This Implicit Baseline champion was selected for $f_1$ at $D=5$ and
$\sigma=0.05$ (experiment 2139, iteration 9). Its current independent
benchmark records eight finite returns and twelve failures across
20 held-out trials. No returned solution meets $\Delta f\leq10^{-8}$;
five meet $\Delta f\leq10^{-2}$. The selection rule, empty-population
failure, and evaluation-counter discrepancy are examined in
Section~\ref{subsec:generated_ses_case}.

\begin{lstlisting}[style=championprogram,caption={Complete generated \texttt{ImprovedSES} champion used in the instance-dependent failure case study.},label={lst:generated_ses_complete}]
import numpy as np

class ImprovedSES:
    def __init__(self):
        self.best_y = float('inf')
        self.evaluations = 0
        self.fitness_history = []

    def __call__(self, problem, budget):
        lb = np.asarray(getattr(problem, 'lower_bound', -5.0), dtype=float)
        ub = np.asarray(getattr(problem, 'upper_bound', 5.0), dtype=float)
        dim = int(getattr(problem, 'dim', len(lb) if hasattr(lb, '__len__') else 3))

        pop_size = 10 * dim  # Empirical choice for population size
        pop = np.random.uniform(lb, ub, size=(pop_size, dim))

        best_x = None
        self.evaluations = 0

        while self.evaluations < budget:
            fitnesses = np.array([problem(x) for x in pop])
            self.evaluations += pop_size
            self.fitness_history.append(np.min(fitnesses))

            if len(self.fitness_history) > 10 and all(f < self.best_y * 1.1 for f in self.fitness_history[-10:]):
                mutation_strength = 0.1 * (ub - lb) * (1 - self.evaluations / budget) * 0.5  # Decrease strength and make it more conservative
            else:
                mutation_strength = 0.1 * (ub - lb) * (1 - self.evaluations / budget)

            idx = np.argsort(fitnesses)[:int(pop_size // 2)]
            parents = pop[idx]
            good_candidates_index = fitnesses[idx] < self.best_y * 1.1
            parents = parents[good_candidates_index]

            mutants = np.clip(parents + np.random.normal(0, mutation_strength, size=parents.shape), lb, ub)
            pop = np.vstack([parents, mutants])

            new_fitnesses = np.array([problem(x) for x in pop])
            self.evaluations += len(new_fitnesses) - len(fitnesses)  # Update evaluation count

            current_best_idx = np.argmin(new_fitnesses)
            if new_fitnesses[current_best_idx] < self.best_y:
                best_x = pop[current_best_idx].copy()
                self.best_y = new_fitnesses[current_best_idx]

        return best_x, float(self.best_y)
\end{lstlisting}

\section{Supplementary SES Diagnostics}
\label{app:ses_diagnostics}

These archived observations and bounded execution checks explain the
offset-sensitive selection rule. They use the historical schedule of
instances 1--20, which had nine finite returns and eleven failures.
They are separate from the current 20 held-out trials reported in
Chapter~\ref{chap:experiments}, which exclude synthesis instance 1 and
have eight finite returns and twelve failures. A finite best-queried
gap in the table does not imply that the program returned a valid point.

\begin{table}[htbp]
    \centering
    \caption{Representative archived trials of \texttt{ImprovedSES} on
    $f_1$, $D=5$, $\sigma=0.05$. Query gaps are obtained from the IOH
    histories; returned gaps and actual evaluation counts are terminal
    measurements. Instance 1 belongs to the historical diagnostic, not
    the current instances 2--21 assessment. The infinite returned gap
    denotes execution failure.}
    \label{tab:ses_instance_diagnostics}
    \small
    \setlength{\tabcolsep}{4pt}
    \begin{tabular}{rrrrrl}
        \hline
        Instance & $f^*$ & Best queried gap & Returned gap & Evaluations & Outcome \\
        \hline
        1 & $79.480$ & $0.007006$ & $0.007006$ & 50,000 & Finite return \\
        3 & $-247.110$ & $5.783063$ & $\infty$ & 150 & Failed \\
        7 & $-1000.000$ & $10.411892$ & $\infty$ & 150 & Failed \\
        16 & $6.480$ & $0.006523$ & $0.006523$ & 50,000 & Finite return \\
        \hline
    \end{tabular}
\end{table}

\FloatBarrier

An additional diagnostic executed the unchanged source with the historical
instance 1--20 schedule and its recorded seeds, stopping surviving searches
after at most 200 objective calls. It reproduced the empty-population \texttt{ValueError}
at 150 calls on all eleven previously failed instances. For instance 3,
the stored best observation after the first pass was approximately
$-241.626$, giving a cutoff of $-265.789$ and zero surviving parents.
For instance 1, the corresponding observation was $87.588$, its cutoff
was $96.347$, and five parents survived. These exception details and local-state values come from the additional
diagnostic, not the archived terminal records.

To isolate the role of the offset, the same diagnostic added $326.59$ to
the observations on instance 3, retaining the initial population and
optimizer/noise seeds. The first 150 queried coordinates and underlying
noise observations were identical; only the returned scalar offset was
changed before selection. Its shifted optimum became $79.48$, and the search continued
to the 200-call diagnostic limit. Conversely, subtracting $326.59$ from
instance 1 made its optimum $-247.11$ and reproduced the empty-population
failure at 150 calls. These bounded checks identify offset sensitivity in
the early selection logic; they do not establish final convergence of the
shifted runs or replace the archived benchmark measurements.
